% !PRINTED 151
\weberpage{151}
%<<CONT
die Co\"efficienten unabh\"angige Variable waren; sie k\"onnen aber
nicht falsch werden, wenn daf\"ur bestimmte numerische Werthe
gesetzt werden, da wenigstens bei den ganzen symmetrischen
Functionen keine Nenner vorkommen, die verschwinden k\"onnen.

Wir machen zun\"achst einige Anwendungen auf besondere
symmetrische Functionen, die f\"ur die Folge wichtig sind. Schon
im \S.\,19 ist bemerkt worden, dass das Differenzenproduct

\weq{1}{\begin{aligned}
P = {}&(\alpha_1 - \alpha_2)\,(\alpha_1 - \alpha_3)\ldots(\alpha_1 - \alpha_n)\\
&\phantom{(\alpha_1 - \alpha_2)\,}(\alpha_2 - \alpha_3)\ldots(\alpha_2 - \alpha_n)\\
&\cdot\;\;\cdot\;\;\cdot\;\;\cdot\;\;\cdot\;\;\cdot\;\;\cdot\;\;\cdot\;\;\cdot\;\;\cdot\;\;\cdot\;\;\cdot\\
&\phantom{(\alpha_1 - \alpha_2)\,(\alpha_1 - \alpha_3)\ldots}(\alpha_{n-1} - \alpha_n)
\end{aligned}}

durch Umstellung zweier Elemente nur das Vorzeichen \"andert.
Das Quadrat dieses Productes wird also bei allen Vertauschungen
zweier der Variablen $\alpha$ und demnach auch bei allen Permu\-tationen
unge\"andert bleiben, d. h. eine symmetrische Function
der Variablen sein. Die Ordnung dieser symmetrischen Function
ist, wie man aus (1) sieht, gleich

\[
(2\,n - 2,\ \ 2\,n - 4 \ldots 2,\, 0).
\]

Wenn wir also die Function $f(x)$, deren Wurzeln die $\alpha_1, \alpha_2 \ldots \alpha_n$
sind, in der Form annehmen

\[
f(x) = a_0\,x^n + a_1\,x^{n-1} + a_2\,x^{n-2} + \cdots + a_{n-1}\,x + a_n,
\]

so dass die h\"ochste Potenz von $x$ nicht gerade den Co\"efficienten
1 hat, so ist $P^2$ eine ganze rationale Function von

\weq{2}{\frac{a_1}{a_0},\ \ \frac{a_2}{a_0} \ldots \frac{a_n}{a_0},}

die durch Multiplication mit $a_0^{2n-2}$ in eine \gesperrt{homogene} Function
von $a_0, a_1, \ldots a_n$ \"ubergeht. Wir nennen diese homogene Func\-tion,
die also durch

\weq{3}{a_0^{2\,n-2}\,P^2 = D}

definirt ist, die \gesperrt{Discriminante der Function} $f(x)$. Setzen
wir $a_0 = 1$, so erhalten wir die \gesperrt{unhomogene Form der
Discriminante}\footnote{Die Discriminante ist, vom Factor $(-\,1)^{\frac{n\,(n-1)}{2}}\ a_0^{2\,n-2}$ abgesehen,
dasselbe, was \gesperrt{Gauss} die Determinante der Function nennt (l. c., art. 6).}.
